%! TEX root = ../../main.tex

%% illustration of relations between hypothesis classes. 

\makeatletter
\newif\iftikz@shading@path

\tikzset{
	% There are three circumstances in which the fading sep is needed:
	% 1. Arrows which do not update the bounding box (which is most of them).
	% 2. Line caps/joins and mitres that extend outside the natural bounding 
	%    box of the path (these are not calculated by PGF).
	% 3. Other reasons that haven't been anticipated.
	fading xsep/.store in=\pgfpathfadingxsep,
	fading ysep/.store in=\pgfpathfadingysep,
	fading sep/.style={fading xsep=#1, fading ysep=#1},
	fading sep=0.0cm,
	shading path/.code={%
			% Prevent this stuff happning recursively.
			\iftikz@shading@path%
			\else\tikz@shading@pathtrue%%
				% \tikz@addmode installs the `modes' (e.g., fill, draw, shade) 
				% to be applied to the path. It isn't usualy for doing more
				% changes to the path's construction.
				\tikz@addmode{%
					\pgfgetpath\pgf@currentfadingpath%
					% Get the boudning box of the current path size including the fading sep
					\pgfextract@process\pgf@fadingpath@southwest{\pgfpointadd{\pgfqpoint{\pgf@pathminx}{\pgf@pathminy}}%
						{\pgfpoint{-\pgfpathfadingxsep}{-\pgfpathfadingysep}}}%%
					\pgfextract@process\pgf@fadingpath@northeast{\pgfpointadd{\pgfqpoint{\pgf@pathmaxx}{\pgf@pathmaxy}}%
						{\pgfpoint{\pgfpathfadingxsep}{\pgfpathfadingysep}}}%
					% Clear the path
					\pgfsetpath\pgfutil@empty%                          
					% Interrupt the path and picture to create a fading.
					\pgfinterruptpath%
					\pgfinterruptpicture%
					\begin{tikzfadingfrompicture}[name=.]
						\path [shade=none,fill=none, #1] \pgfextra{%
							% Set the softpath. Any transformations in #1 will have no effect.
							% This will *not* update the bounding box...
							\pgfsetpath\pgf@currentfadingpath%
							% ...so it is done manually.
							\pgf@fadingpath@southwest
							\expandafter\pgf@protocolsizes{\the\pgf@x}{\the\pgf@y}%
							\pgf@fadingpath@northeast%
							\expandafter\pgf@protocolsizes{\the\pgf@x}{\the\pgf@y}%
						};
						% Now get the bounding of the picture.
						\xdef\pgf@fadingboundingbox@southwest{\noexpand\pgfqpoint{\the\pgf@picminx}{\the\pgf@picminy}}%
						\xdef\pgf@fadingboundingbox@northeast{\noexpand\pgfqpoint{\the\pgf@picmaxx}{\the\pgf@picmaxy}}%
						%
					\end{tikzfadingfrompicture}%
					\endpgfinterruptpicture%
					\endpgfinterruptpath%
					% Install a rectangle that covers the shaded/faded path picture.                                
					\pgfpathrectanglecorners{\pgf@fadingboundingbox@southwest}{\pgf@fadingboundingbox@northeast}%
					% Make the fading happen.
					\def\tikz@path@fading{.}%
					\tikz@mode@fade@pathtrue%
					\tikz@fade@adjustfalse%10pt
					% Shift the fading to the mid point of the rectangle
					\pgfpointscale{0.5}{\pgfpointadd{\pgf@fadingboundingbox@southwest}{\pgf@fadingboundingbox@northeast}}%
					\edef\tikz@fade@transform{shift={(\the\pgf@x,\the\pgf@y)}}%
				}%
			\fi%
		}
}

\begin{tikzpicture}[scale=1,
	]
	\begin{axis}[width=1.15\linewidth, height=2.5cm,
			axis lines=none,  % don't print axis lines
			yticklabels={,,}, xticklabels={,,},
			xmin=-10, xmax=10, x axis line style={-},
			ymin=-1, ymax=2, y axis line style={-},
		]
		%\draw [<->, solid, draw=orange, line width = 1mm] (axis cs:-10,0) -- (axis cs:10,0);
		\node[align=center] at (axis cs:-8.5,1) {\small Easy};
		\node[align=center] at (axis cs:0,1) {\small Hard};
		\node[align=center] at (axis cs:7.5,1) {\small Even Harder};
		\path [left color=blue, right color=orange,
			shading path={draw=transparent!0, line width = 1mm, <->}]
		(axis cs:-10,0) -- (axis cs:10,0);

	\end{axis}

\end{tikzpicture}%
